\section{Sonuç}\label{sec:sonuc}

Hızlı bir ILS'te double-bridge vuruşlarının nereye yönlendirileceğini, öğrenilmiş ve eğitimsiz
sinyalleri ön-kayıtlı bir doğrulayıcı testte karşılaştırarak inceledik. Uygulayıcı için sonuç
basittir: vuruşları en uzun tur kenarlarına (ön-kayıtlı vekilde uç noktaların yerel ölçeğine göre
düzeltilmiş) yönlendirmek, bu motor ve 10~sn bütçede rastgele vuruştan geliştirme ve doğrulama
setlerinde tutarlı biçimde iyidir; sınanan öğrenilmiş skorlar karşısında bunun için model eğitmek
gerekmedi. Bu eğitimsiz vekil, LKH-3 turlarıyla eğitilmiş GNN'ye ve kazanç hedefli MLP'ye karşı
kötü-değil testlerini geçti, doğrulama setinde ortalama gap'i \%1,094'ten \%0,922'ye indirdi ve
keşifsel karşılaştırmalarda öğrenilmiş skorların önünde kaldı.
Bu başarım üç tamamlayıcı özellikle tutarlıdır: kazanca hizalı hedefleme, ihmal edilebilir
hesaplama maliyeti ve sürekli güncel havuz yapısı. Keşifsel bir kolda yerel ölçek düzeltmesi olmadan en
uzun kenarları hedeflemek de rastgeleye karşı benzer bir kazanç sağladı ve vekilden ortalamada
ayrışmadı ($+0{,}009$ puan); bu, yararın bu motorda doğrudan kenar uzunluğu bilgisinden geldiğini
düşündürür ve en basit vekile karşı sınama ilkesinin kendi önerimize de yararla uygulanabildiğini
gösterir.
Koşulardan önce kaydedilen keşifsel bloklar bu tabloyu genişletir: hedefli global vuruş, klasik
rastgele segment vuruşundan 1,04 puan iyidir; vekilin rastgeleye ve GNN'ye karşı avantajı 60~sn'de
daralarak sürer; geliştirmede seçilen havuz kuralının çevresindeki plato holdout2'de de görülür.
Hakem odaklı bloklar tabloyu netleştirir. Yayımlanmış NeuroLKH modeli, süresi bütçeden düşüldüğünde
vekilden ayrışmaz; süresi sayılmadığında problemlerin çoğunda küçük bir farkla öndedir. Vuruş sayısı
eşitlendiğinde öğrenilmiş skorlar vekilden kötü değildir ve vuruş başına en iyi sinyal eğitimsiz kNN
sırasıdır. Yani eşit süredeki sıralamayı sinyalin niteliğinden çok maliyeti belirler; vuruş başına en iyi
sinyal de, model süresi sayıldığında eşit sürede en düşük ortalama gap'i veren sinyal de eğitimsizdir. LKH-3'te ise hedefleme ölçülebilir bir kazanç getirmedi; bulgular hızlı,
2-opt/Or-opt tabanlı ILS için geçerlidir.
% kaynak: E608, E610, E965, E972, E1047, E1077, E1078, E1091, E1106, E1107, E1128-E1131, E1172

Koşu içi çevrimiçi öğrenme ise işaretli ödülle kendi ödülünü artırırken son kaliteyi düşürdü: vuruş
başına ortalama ödül, yerel aramanın onardığı nötr vuruşları ödüllendirir. Kredi kabul edilen
iyileşmelere dayandırıldığında bu zarar kayboldu, ama sabit modele bir kazanç da eklenmedi. Seyrek
iyileşmeli aramalarda kredi tasarımı bu nedenle öncelikli bir sorudur.
% kaynak: E1150-E1152 Yöntemsel ders de aynı yöndedir: öğrenilmiş bir arama bileşeni
rastgele tabana değil, en basit kazanca hizalı eğitimsiz vekile karşı ve eşit duvar-saatinde
değerlendirilmelidir.

Sonraki adımlar, karşılaştırmanın dakikalar mertebesini aşan bütçelerde tekrarı; vekil, kNN sırası
ve NeuroLKH arasındaki farkların ön-kayıtlı bir doğrulayıcı testle sınanması; ham ve ölçekle
düzeltilmiş uzunluğun ön-kayıtlı bir eşdeğerlik testiyle karşılaştırılması ve LK düzeyindeki
aramalarda hedeflemenin neden etkisiz kaldığının incelenmesidir.
